Today it is common to build systems where several LLM agents work together. One agent browses the web, another reads files, a third keeps notes in memory, and a fourth sends emails or makes payments. All of them read text that nobody has checked. If an attacker hides a few lines inside a web page or a document, an agent may follow those lines instead of the user's request. This is called indirect prompt injection. The version we worry about most is split into pieces: an address hidden in one document, a note saved to memory, and much later an email that uses the address. Each message looks harmless on its own, so a guardrail that checks one message at a time cannot see the attack. The proof is only in the history of the session.

In this project we built MAPIS, a monitor that sits on the communication layer between the agents and remembers the session. Every message, tool call, tool output and memory access goes through a FastAPI tap. MAPIS keeps a small record of the user's goal, where each address, link or id first appeared, any claims of authority, and how each agent normally behaves. A fine-tuned DeBERTa-v3-small model reads the current event together with this record and four signals (instructions inside data, drift from the goal, cross-hop provenance and unusual behaviour) and gives a trust score between $0$ (malicious) and $1$ (safe). Depending on the score, the message is passed, flagged, quarantined for an operator, or blocked, and the operator sees a trace of where the attack came from. The system runs on both LangGraph and AutoGen.

We also built MAPIS-Bench, six test sets that include multi-hop attacks paired with harmless ``twin'' sessions and an attacker who tries to hide from MAPIS. On 2,304 held-out sessions, averaged over five training runs, MAPIS caught $98.0\pm0.6\%$ of attacks with a false-positive rate of $0.9\pm0.8\%$ and $98.5\%$ accuracy, at 63 ms per event. The same model without session memory needed a $33.2\%$ false-positive rate to reach $94.1\%$ recall, and when we limited it to $5\%$ false alarms it caught only $4.3\%$ of multi-hop attacks, against $100\%$ for MAPIS. We compared against Llama Guard, Prompt Guard 2, ProtectAI, NeMo Guardrails and an LLM judge on the same sessions. With a real LLM agent on AgentDojo, MAPIS brought the attack success rate down from $15.3\%$ to $0\%$. All four project targets were met. Two harder test sets showed one clear gap: attacks that only change a value, add a privilege flag or pass a destination through another agent. After we added three checks on action arguments, without retraining, recall on a set written by our own team went from $40\%$ to $95\%$ and on attacks against the session record from $55.7\%$ to $80.4\%$, with no new false positives. Since we designed these checks after seeing the team set's misses, a fresh set is still needed to confirm that result.
